\documentclass[12pt,a4paper]{article}

%!TEX program = pdflatex

% =============================================================================
% 1. PHÔNG CHỮ & NGÔN NGỮ (Hỗ trợ Tiếng Việt)
% =============================================================================
\usepackage[utf8]{inputenc}
\usepackage[T5]{fontenc}
\usepackage[vietnamese]{babel}

% =============================================================================
% 2. BỐ CỤC TRANG (LAYOUT)
% =============================================================================
\usepackage[margin=2.5cm]{geometry}
\usepackage{setspace}
\onehalfspacing % Khoảng cách dòng 1.5 (tùy chọn)

% =============================================================================
% 3. ĐỒ HỌA & BIỂU ĐỒ (TikZ & DFD Styles)
% =============================================================================
\usepackage{graphicx}
\usepackage{tikz}
\usetikzlibrary{arrows.meta, positioning, shapes.geometric}

% Định nghĩa style cho DFD (Data Flow Diagram)
\tikzset{
    process/.style={ellipse, draw, minimum width=3cm, minimum height=1.5cm, align=center},
    entity/.style={rectangle, draw, minimum width=2.5cm, minimum height=1cm, align=center},
    datastore/.style={
        draw=none, align=center, minimum width=3cm,
        append after command={
            [shorten <= -0.2cm, shorten >= -0.2cm]
            (\tikzlastnode.north west) edge (\tikzlastnode.north east)
            (\tikzlastnode.south west) edge (\tikzlastnode.south east)
        }
    },
    arrow/.style={-Stealth, thick}
}

% =============================================================================
% 4. BẢNG BIỂU & DANH SÁCH
% =============================================================================
\usepackage{array}
\usepackage{booktabs}
\usepackage{multirow}
\usepackage{longtable}
\usepackage{float}
\usepackage{enumitem}
\usepackage{tabularx}
\usepackage{makecell}
\usepackage{colortbl}

% =============================================================================
% 5. TIỆN ÍCH VĂN BẢN & LIÊN KẾT
% =============================================================================
\usepackage{xcolor}
\usepackage{titlesec}
\usepackage{tocloft}
\usepackage{xurl}
\usepackage{seqsplit}
\usepackage{hyperref} % Nên để hyperref ở cuối nhóm package

\hypersetup{
    colorlinks=false,
    hidelinks,
    pdftitle={Báo cáo Kỹ thuật},
    pdfauthor={Tên của bạn}
}

% =============================================================================
% 6. ĐỊNH DẠNG ĐOẠN VĂN & ĐÁNH SỐ MỤC
% =============================================================================
\setlength{\parindent}{0pt}
\setlength{\parskip}{6pt}

\setcounter{secnumdepth}{3} % Đánh số đến mức subsubsubsection

% Định dạng hiển thị số: I, 1, 1.1
\renewcommand{\thesection}{\Roman{section}}            
\renewcommand{\thesubsection}{\arabic{subsection}}     
\renewcommand{\thesubsubsection}{\thesubsection.\arabic{subsubsection}} 

% Thiết kế tiêu đề (Titleformat)
\titleformat{\section}{\large\bfseries}{\thesection.}{0.5em}{}
\titleformat{\subsection}{\normalsize\bfseries}{\thesubsection.}{0.5em}{}
\titleformat{\subsubsection}{\normalsize\itshape\bfseries}{\thesubsubsection.}{0.5em}{}

% Cấu hình Paragraph làm subsubsubsection (Mức 4)
\titleformat{\paragraph}[block]{\normalsize\bfseries}{}{0pt}{}
\titlespacing*{\paragraph}{0pt}{6pt}{0pt}

% =============================================================================
% 7. MỤC LỤC (TOC)
% =============================================================================
\renewcommand{\cftsecaftersnum}{.}
\renewcommand{\cftsubsecaftersnum}{.}
\renewcommand{\cftsubsubsecaftersnum}{.}

% =============================================================================
% 8. CUSTOM COMMANDS & COLUMN TYPES
% =============================================================================
\newcolumntype{C}[1]{>{\centering\arraybackslash}p{#1}}
\newcolumntype{L}[1]{>{\raggedright\arraybackslash}p{#1}}
\newcolumntype{Y}{>{\raggedright\arraybackslash}X}

\newcommand{\subsubsubsection}[1]{\paragraph{#1}}
\newcommand{\tablesetup}{\small\setlength{\tabcolsep}{3pt}\renewcommand{\arraystretch}{1.2}}
\newcommand{\code}[1]{{\ttfamily\seqsplit{#1}}}


% =================================
% Dùng cho phần Sequence Diagram
\newcommand{\seqpdf}[1]{%
\IfFileExists{#1}{%
    \includegraphics[page=1,width=1.05\textwidth,height=0.95\textheight,keepaspectratio]{#1}%
}{%
    \fbox{%
        \parbox[c][0.28\textheight][c]{0.9\textwidth}{%
            \centering
            Chưa tìm thấy file \texttt{#1}.\\[4pt]
            Hãy upload đúng file PDF này lên Overleaf.
        }%
    }%
}%
}
